\subsection{Related Work}
\label{vfree:sec:related}
% We discuss key works related to {\vfreename}.
% \vfreebrand is the first work that studies IO memory protection in virtualized environments with nested transaction, on modern server hardware architectures.

\para{Bare-metal IO memory protection.}
There is a large and active body of work on enabling IO memory protection with near-zero overheads for the bare metal case~\cite{ben2007price, Malka2015rIOMMU, malka2015efficient, peleg2015utilizing, Markuze2016true, markuze2018damn, RajapakshaIOMMU, Rubin2024FastSafe}. There is also work on rearchitecting IOMMU (\textit{e.g.}, replacing page tables with circular tables~\cite{Malka2015rIOMMU}
    and using hardware support for security checks~\cite{RajapakshaIOMMU}). Collectively, these works have led to mechanisms that enable achieving strict IO memory protection with minimal overheads. 

{\vfreename} builds upon these works for the bare metal case. For example, persistent
mapping for receive buffers has been adopted by Linux (Rx datapath optimizations in \S\ref{vfree:sec:background-datapath}); similarly, as discussed in \S\ref{vfree:sec:impl}, we use the techniques developed in~\cite{Rubin2024FastSafe}. As shown in Figure~\ref{vfree:fig:ablation}, these mechanisms are helpful but not sufficient in virtualized environments---the key focus of our work. Indeed, we show in \S\ref{vfree:sec:inefficiencies} that the performance bottlenecks are quite different in virtualized environments and in fact may require design choices that make very different tradeoffs (\textit{e.g.}, {\vfreename} reduces the number of VM exits by increasing the address translation cost).
    
% % Prior work studies the overhead of IOMMU protection in bare-mental environments.
% Markuze et al.\ proposed shadow DMA buffers that restrict devices to
% permanently mapped memory regions, eliminating IOTLB invalidations and
% achieving up to 5$\times$ higher throughput than strict mode~\cite{Markuze2016true}.  
% DAMN specializes this approach for networking by allocating packet buffers from a
% permanently mapped pool~\cite{markuze2018damn}. Persistent
% mapping for receive buffers has been adopted by Linux (Rx datapath optimizations in \S\ref{vfree:sec:background-datapath}).
% % subsystem, so the copy-based approach cannot support zero-copy workloads.
% % and DAMN is inapplicable to non-networking devices such as NVMe.  
% In fact, the CPU and memory bandwidth cost of copying scales with IO throughput, 
%     limiting the approach at higher link speeds.
% F\&S exploits IO PWCs to reduce the cost of IOTLB misses, 
%     achieving strict safety with near-zero overhead on bare metal~\cite{Rubin2024FastSafe}. 
% F\&S's IOVA allocation strategy and page pools apply to virtualized environments, 
%     but are insufficient:
%     under nested translation, the dominant overhead shifts from the cost of IOTLB misses 
%     to VM exits. \vfreebrand builds on the bare-metal foundations and reduces VM exits due to 
%     cache invalidations, the dominant bottleneck as link speeds increase.

% \vfreebrand requires no IOMMU hardware modification and is 
%     complementary to research that proposes to rearchitect
%     IOMMU (e.g., replacing page tables with circular tables~\cite{Malka2015rIOMMU}
%     and hardware support for security checks~\cite{RajapakshaIOMMU}).
% Faster IOMMU hardware would further amplify gains of \vfreebrand.
% Several proposals rearchitect IOMMU hardware to reduce translation or
% invalidation costs.  Malka et al.\ (rIOMMU) replace the page-table hierarchy
% with a flat circular table for ring-buffer devices~\cite{Malka2015rIOMMU}.  
% Rajapaksha et al.\ proposed a hardware indicator-bit to close the exploitable stale-entry window created by deferred invalidation~\cite{RajapakshaIOMMU}.  
% These approaches are complementary to \vfreebrand: reducing the hardware 
% cost of each invalidation could further amplify the gains from \vfreebrand's software-side batching and serialization elimination.

\para{vIOMMU emulation.}
Amit et al.\ presented an vIOMMU emulation for
    unmodified guests~\cite{Amit2011vIOMMU}. Prior work also optimized other virtualization overhead including 
    memory pinning~\cite{Tian2020coIOMMU} and interrupts~\cite{Gordon2012eli, Har2013paraIO}. All the above work was conducted at 10--100\,Gbps. \vfreebrand extends these works in two ways. First, it presents new insights related to performance bottleneck at 400\,Gbps (\S\ref{vfree:sec:inefficiencies}); and second, {\vfreename} addresses these bottlenecks with software techniques that require no paravirtualization or hardware modifications.

%     with 
%     two optimizations: sidecore offloading and optimistic teardown. 
% Sidecore offloading requires hypervisor modifications and a dedicated core; 
% \vfreebrand uses a push-based design where vCPUs enqueue invalidations into per-CPU buffers without waiting
% or dedicating a core. Optimistic teardown recovers line rate at 10\,GbE,
%     but weakens security guarantees; \vfreebrand preserves strict safety.

\if 0 % this conference is way too weird
LA-vIOMMU added a fixed IOTLB for long-lived mappings and a guest-side DMA mapping cache 
    to reduce VM-exit frequency~\cite{Lv2022LAvIOMMU}.
LA-vIOMMU's guest-side cache optimization targets the map/unmap exit cost; 
    under nested translation, mapping overhead is already low. 
LA-vIOMMU also requires hardware modifications to the IOTLB.
\fi 


% coIOMMU introduced cooperative DMA buffer tracking between
%guest and hypervisor for fine-grained memory pinning~\cite{Tian2020coIOMMU}. 
% coIOMMU targets the memory pinning problem --- determining which guest pages to pin in host memory---which 
% is orthogonal to the IOTLB invalidation serialization that \vfreebrand addresses.  

% Parallel efforts have targeted other virtualization overheads.  Gordon et al.\
% (ELI) and Har'El et al.\ reduced interrupt and exit costs for assigned
% devices, recovering near-native throughput~\cite{Gordon2012eli, Har2013paraIO}.  
% These optimizations are complementary to \vfreebrand, which focuses on the IOMMU emulation path rather than the interrupt delivery path.




\para{virtIO IOMMU.} \vfreebrand targets hardware-assisted vIOMMU, not virtio-IOMMU;
    despite that, the {\vfreename} ideas may apply.
vIOMMU offers VMs with native hardware interfaces, 
    ensuring compatibility with nested translation and modern drivers.
It is supported by mainstream systems such as QEMU~\cite{qemu_10_0}.
On our platform (\S\ref{vfree:sec:inefficiencies}),
    virtio-IOMMU can only hit 51\% of the vanilla 
    vIOMMU throughput with nested translation.
%    , and just 9\% compared to vIOMMU Off and \vfreename{}.}

% \tianyin{@Siyuan: do we have one result that virtio-IOMMU is slower?}



% Besides hardware-assisted vIOMMU, IOMMU can be virtualized through virtio-IOMMU, a paravirtualization method 
%    that requires a synchronized software stack involving both a guest-side frontend driver and a host-side backend.
% Each DMA operation requires a hypercall, mutex acquisition, and software tree traversal.
% \schaic{I replaced VM exit with hypercall to reinforce their VM exit nature. @Owen Is it one hypercall on both map and unmap time? If so, we should explicitly state it. Also, let's explain from functionality wise what is the hypercall doing.}
% Although this implementation is simple, it is considerably less efficient compared to Intel VT-d,
% which does not require a VM exit when new mappings are built.
% \schaic{@Owen, let's include a percentage number based on our profiling. aka vanilla virtio vs. vanilla nested strict.}

% The virtio-vIOMMU also introduces a non-architectural software abstraction that precludes the use of unmodified guest vendor drivers.
% These high-performance proprietary stacks and userspace drivers (e.g., DPDK) often bypass standard kernel DMA APIs to interact directly with architectural register sets.
% Because virtio-IOMMU exposes a generic software device rather than the expected hardware interfaces, 
%    these drivers may fail to detect a compatible IOMMU, breaking support for advanced features like p2p DMA and PASID.
% \schaic{@Owen, can you help search for reference to support our argmument.s}
    
% Building upon the vIOMMU framework, our work provides virtual machines with the exact same architectural interface used by a bare-metal kernel. 
% It ensures compatibility with emerging hardware-assisted nested translation and modern drivers.
% Our approach targets the production-standard interface used in modern data centers:
%    \vfreebrand eliminates vIOMMU overhead while maintaining strict safety protection.  
    % allowing for optimizations in DMA mapping and IOTLB invalidation that remain entirely transparent to the guest while utilizing hardware-assisted nested translation.
